% XRPL Lending Protocol Hackathon · CY-HACK · Developer Feedback (3 pages)
% Build : tectonic -X compile deck/feedback-report.tex --outdir deck
% Captures : node deck/capture-explorer.mjs   (→ deck/shots/*.png)
% DA du deck : navy 001C5C · bleu 006AFF · ciel 6DC3FF · corps 5F666E
\documentclass[a4paper,10pt]{extarticle}

\usepackage[margin=11mm,top=9.5mm,bottom=11mm,footskip=5.5mm]{geometry}
\usepackage{fontspec}
\usepackage[table]{xcolor}
\usepackage{tabularx,booktabs,array,colortbl}
\usepackage{multicol}
\usepackage{graphicx}
\usepackage{xfp}
\usepackage{tcolorbox}
\usepackage{enumitem}
\usepackage{microtype}
\usepackage{fancyhdr}
\usepackage{ragged2e}
\usepackage{lastpage}
\tcbuselibrary{skins}
\graphicspath{{deck/shots/}{shots/}}

% ── DA ──
\definecolor{navy}{HTML}{001C5C}
\definecolor{blue}{HTML}{006AFF}
\definecolor{sky}{HTML}{6DC3FF}
\definecolor{body}{HTML}{5F666E}
\definecolor{muted}{HTML}{8A9096}
\definecolor{card}{HTML}{F6F7F8}
\definecolor{line}{HTML}{E2E5E8}
\definecolor{bluecard}{HTML}{EBF4FF}
\definecolor{chipbg}{HTML}{BAEAFF}
\definecolor{chiptx}{HTML}{0045C6}
\definecolor{peachbg}{HTML}{FAEDE7}
\definecolor{peachfg}{HTML}{C4562F}

% Faces déclarées explicitement : laissée à la résolution automatique,
% HelveticaNeue.ttc sort \bfseries en Bold Italic. Pas d'option Color= non plus :
% elle casse xdvipdfmx avec tcolorbox (« Invalid object type »).
\setmainfont{Helvetica Neue}[
  UprightFont = {Helvetica Neue}, BoldFont = {Helvetica Neue Bold},
  ItalicFont = {Helvetica Neue Italic}, BoldItalicFont = {Helvetica Neue Bold Italic}]
\setmonofont{Menlo}[Scale=0.84,
  UprightFont = {Menlo Regular}, BoldFont = {Menlo Bold},
  ItalicFont = {Menlo Italic}, BoldItalicFont = {Menlo Bold Italic}]
\newfontfamily\hdfont{Poppins}[BoldFont={Poppins SemiBold}]
\makeatletter\renewcommand\normalsize{\@setfontsize\normalsize{9.6}{12.1}}\makeatother
\AtBeginDocument{\color{body}\normalsize}

\setlength{\parindent}{0pt}
\setlength{\parskip}{3.4pt}
\linespread{1.02}
\raggedcolumns
% Les listes posent une pénalité négative (−51) avant et après elles : avec des
% colonnes non étirées, c'est la coupure la moins chère, et multicols coupe là en
% laissant un vide. On la neutralise, et on interdit les lignes orphelines ou veuves
% (LaTeX recharge \clubpenalty depuis \@clubpenalty à chaque paragraphe).
\makeatletter
\@beginparpenalty=0 \@endparpenalty=0 \@itempenalty=0
\@clubpenalty=10000 \clubpenalty=10000 \widowpenalty=10000
\makeatother
\setlength{\columnsep}{6mm}
\setlength{\columnseprule}{0.4pt}
\renewcommand{\columnseprulecolor}{\color{line}}
\arrayrulecolor{line}
\renewcommand{\tabularxcolumn}[1]{m{#1}}
\renewcommand{\arraystretch}{1.28}

\pagestyle{fancy}\fancyhf{}
\renewcommand{\headrulewidth}{0pt}
\renewcommand{\footrulewidth}{0.4pt}
\fancyfoot[L]{\scriptsize\color{muted}CY-HACK · XRPL Lending Protocol Hackathon}
\fancyfoot[R]{\scriptsize\color{muted}\thepage\,/\,\pageref{LastPage}}

% ── briques ──
\newcommand{\hh}[1]{\par\addvspace{8pt}{\raggedright\hyphenpenalty=10000\exhyphenpenalty=10000\hdfont\color{navy}\fontsize{10.4}{12.2}\selectfont\bfseries #1\par}\nobreak\addvspace{2.5pt}}
\newcommand{\hs}[1]{{\color{navy}\bfseries #1}}
\newcommand{\tm}[1]{{\ttfamily\color{navy}#1}}
\newcommand{\lab}[1]{{\color{blue}\bfseries #1}\enspace}
% hash abrégé 8+8 dans les tableaux ; le hash complet est dans FEEDBACK.md
\newcommand{\hsh}[2]{{\ttfamily\fontsize{6}{7}\selectfont\color{body}#1…#2}}
% hashes sous une preuve : jamais séparés de ce qui précède, jamais coupés entre deux colonnes
\newcommand{\ev}[1]{\par\nobreak\noindent\begin{minipage}{\linewidth}\raggedright\ttfamily\fontsize{5.3}{6.6}\selectfont\color{muted}#1\end{minipage}\par}
\newcommand{\chip}[3]{\tcbox[on line,colback=#1,colframe=#1,boxsep=0pt,left=2.4pt,right=2.4pt,
  top=1pt,bottom=1pt,arc=1.5pt]{\fontsize{5.8}{6}\selectfont\color{#2}\bfseries\MakeUppercase{#3}}}

% ── captures d'écran de l'explorer ──
% Toutes au même grossissement : 460 px CSS (2415 px à ×5,25) remplissent la largeur
% utile de la carte, soit ≈ 0,56 pt par px CSS : le texte de l'explorer sort entre 7 et
% 10 pt, comme les tableaux et le corps. Pas de titre de transaction (≈ 22 pt) : la
% pastille de statut seule (*-badge), le type est dans la légende.
\newtcolorbox{screen}{enhanced,colback=black,colframe=black,boxrule=0pt,arc=2.4pt,
  left=3.5pt,right=3.5pt,top=3pt,bottom=3pt,boxsep=0pt,before skip=4pt,after skip=2pt}
\newcommand{\shot}[1]{\includegraphics[scale=\fpeval{\linewidth/2415bp}]{#1.png}\par}
\newcommand{\shotw}[1]{\includegraphics[width=\linewidth]{#1.png}\par}
\newcommand{\capt}[1]{{\raggedright\fontsize{6.3}{7.6}\selectfont\color{muted}#1\par}\addvspace{2pt}}
% capture + légende + hash : un seul bloc insécable, jamais coupé entre deux colonnes
\newenvironment{shotblock}{\par\addvspace{6pt}\noindent\begin{minipage}{\linewidth}\setlength{\parskip}{0pt}}{\end{minipage}\par\addvspace{5pt}}
\setlist[itemize]{leftmargin=8pt,itemsep=1.7pt,topsep=2.4pt,parsep=0pt,label={\color{blue}\textbullet}}

\begin{document}

% ═══════════ BANDEAU ═══════════
\begin{tcolorbox}[colback=navy,colframe=navy,boxrule=0pt,arc=2.5pt,left=8pt,right=8pt,top=6pt,bottom=6pt,boxsep=0pt]
\begin{minipage}[c]{0.66\linewidth}
{\hdfont\color{white}\fontsize{15}{16.5}\selectfont\bfseries XRPL Lending Protocol · Developer Feedback}\\[2pt]
{\fontsize{7.2}{8.6}\selectfont\color{sky}Hackathon DeVinci Blockchain × Ripple · 12 and 13 September 2026}
\end{minipage}\hfill
\begin{minipage}[c]{0.32\linewidth}\raggedleft
{\hdfont\color{white}\fontsize{9}{10}\selectfont\bfseries CY-HACK}\\[1.5pt]
{\fontsize{6.4}{7.6}\selectfont\color{sky}Track 1 · open-ended vault\\
XLS-65, XLS-66 and TokenEscrow\\
\texttt{github.com/charlyppr/xrpl-lending}}
\end{minipage}
\end{tcolorbox}

\vspace{4pt}
{\setlength{\tabcolsep}{4pt}\fontsize{6.7}{8.2}\selectfont
\begin{tabularx}{\linewidth}{@{}>{\raggedright\arraybackslash}X|>{\raggedright\arraybackslash}X|>{\raggedright\arraybackslash}X|>{\raggedright\arraybackslash}X@{}}
{\color{blue}\bfseries NETWORK} & {\color{blue}\bfseries LIBRARY} & {\color{blue}\bfseries FIRST TRANSACTION} & {\color{blue}\bfseries VERIFICATION}\\
Custom Hackathon Devnet · \tm{rippled} 3.4.0-rc1, hackathon branch · \tm{LendingProtocolV1\_1} and \tm{fixCleanup3\_4\_0} enabled &
\hs{\texttt{xrpl@4.6.0}} · the 15 XLS-65/66 transaction types are typed · no raw JSON &
\hs{50 minutes}, of which 45 on blocked ports (finding 8) &
\hs{143 of 143} cited hashes validated on the ledger · docs checked against XLS-65, XLS-66 and xrpl.org · library claims against the \tm{rippled} and \tm{xrpl.js} source\\
\end{tabularx}}

\vspace{5pt}
{\renewcommand{\arraystretch}{1.2}\setlength{\tabcolsep}{3.5pt}\fontsize{7}{8.4}\selectfont
\begin{tabularx}{\linewidth}{@{}c >{\raggedright\arraybackslash}X l l@{}}
\rowcolor{card}
{\bfseries\color{navy}\#} & {\bfseries\color{navy}Finding} & {\bfseries\color{navy}Category} & {\bfseries\color{navy}Checked against}\\
\midrule
1 & \tm{fixCleanup3\_4\_0} changes impairment timing and the \tm{LoanSet} counterparty signature. xrpl.org does not list it. & docs, amendments & ledger, XLS-66, rippled\\
\rowcolor{card}
2 & A late \tm{LoanPay} without \tm{tfLoanLatePayment} returns \tm{tecEXPIRED}. The xrpl.org \tm{LoanPay} reference does not list this failure. & docs, error codes & ledger, XLS-66, xrpl.org\\
3 & The first-loss capital covered 0.5\,\% of a defaulted loan. The broker then withdrew the remaining cover. & design, naming & ledger, xrpl.org\\
\rowcolor{card}
4 & xrpl.org omits the sole-holder exception to share value. \tm{LossUnrealized} is absent from the node while zero. & docs, UX & ledger, XLS-65, xrpl.org\\
5 & \tm{loan\_info} and \tm{loan\_broker\_info} do not exist. \tm{mpt\_holders} needs Clio, which the devnet does not expose. & missing primitive & ledger, xrpl.org\\
\rowcolor{card}
6 & Four result codes each cover situations that call for different actions. & error codes & ledger\\
7 & The hackathon build reverts the V1.1 \tm{LoanBrokerSet} restriction. The event brief does not mention it. & event docs & ledger, rippled\\
\rowcolor{card}
8 & The devnet listens only on ports 51233 and 51234. A filtered port and a stalled node show the same signature. & infrastructure & network tests\\
9 & Closing a vault to deposits requires setting \tm{AssetsMaximum} to \tm{AssetsTotal}. & missing primitive & ledger, xrpl.org\\
\rowcolor{card}
10 & Vault shares carry \tm{lsfMPTCanTrade}. \tm{OfferCreate} with a share amount returns \tm{temDISABLED}. & docs, MPT & ledger, xrpl.org\\
\end{tabularx}}

\vspace{4pt}
\begin{multicols}{2}\RaggedRight

% ───────────── 1 ─────────────
\hh{1 · \tm{fixCleanup3\_4\_0} is enabled and not listed on xrpl.org}

The \tm{feature} command reports \tm{fixCleanup3\_4\_0} as enabled on the devnet. The xrpl.org
\emph{Known Amendments} page does not list it. The amendment changes two lending behaviours.

\lab{Impairment.}XLS-66, section 3.10.4.2, condition 9: with \tm{fixCleanup3\_4\_0},
\tm{tfLoanImpair} fails while \tm{currentTime <= NextPaymentDueDate}. The xrpl.org tutorial
\emph{Manage a Loan} says a broker can ``\emph{manually impair a loan before a payment due date
passes}''.

{\setlength{\tabcolsep}{3pt}\fontsize{7}{8.4}\selectfont
\begin{tabularx}{\linewidth}{@{}>{\raggedright\arraybackslash}X l@{}}
\rowcolor{card}
{\bfseries\color{navy}Time relative to \texttt{NextPaymentDueDate}} & {\bfseries\color{navy}\texttt{tfLoanImpair}}\\
\midrule
31 s before & \tm{tecTOO\_SOON}\\
\rowcolor{card}
13 s before & \tm{tecTOO\_SOON}\\
5 s after & \tm{tesSUCCESS}\\
\end{tabularx}}
\ev{5 s after: 5486020BA8FF227C58409042F96BDC07246FDACC3F9B0EF24749555E319E972E}

\lab{Counterparty signature.}rippled PR \#8162: with \tm{fixCleanup3\_4\_0},
\tm{CounterpartySignature} uses the \tm{CPT} signing prefix. In \tm{xrpl@4.6.0},
\tm{signLoanSetByCounterparty} signs with \tm{encodeForSigning}, the previous prefix, and has no
option to change it. Its \tm{LoanSet} fails local checks with \tm{Counterparty: Invalid
signature}. The same \tm{LoanSet} signed with \tm{encodeForSigningCounterparty} from
\tm{ripple-binary-codec} returned \tm{tesSUCCESS}.
\ev{23631C3610DA97926E8BF9FC5BC4E76E504C9DFB2683C928E825DE0138397B2B}

% ───────────── 2 ─────────────
\hh{2 · A late \tm{LoanPay} without \tm{tfLoanLatePayment} returns \tm{tecEXPIRED}}

A loan 10 s past \tm{NextPaymentDueDate}, within its grace period, rejected four \tm{LoanPay}
transactions with \tm{tecEXPIRED}: the instalment, the instalment plus \tm{LatePaymentFee}, the
full balance, and a payment with \tm{tfLoanFullPayment}. The same payment with
\tm{tfLoanLatePayment} (262144) returned \tm{tesSUCCESS}.

\begin{itemize}
\item XLS-66, section 3.11.4.2, condition 11 defines this failure. The xrpl.org \tm{LoanPay} reference lists eight error codes, and \tm{tecEXPIRED} is not among them.
\item xrpl.org describes \tm{tfLoanLatePayment} as ``\emph{Indicates that the borrower is making a late loan payment}''.
\item The explorer shows \tm{tfLoanImpair} and \tm{tfLoanDefault} by name, and \tm{tfLoanLatePayment} as \tm{0x00040000}.
\item No field of the \tm{Loan} node changed during 85 s of lateness, and no field holds the accrued late interest. A payment of \tm{ceil(PeriodicPayment) + LoanServiceFee + LatePaymentFee} returned \tm{tecINSUFFICIENT\_PAYMENT}. A payment of 2\,765\,001 drops returned \tm{tesSUCCESS}, and 765\,003 drops were taken.
\end{itemize}

\begin{shotblock}
\begin{screen}
\shot{f1a-badge}\shot{f1a-status}
\end{screen}
\capt{Explorer. \texttt{LoanPay} of 1.030001 XRP without flags.}
\ev{96C3870480E4CB0C9E3B925F2965347DA0815A486CDB4B49545928E32D6B0BA4}
\end{shotblock}
\begin{shotblock}
\begin{screen}
\shot{f1b-badge}\shot{f1b-flags}
\end{screen}
\capt{Explorer. \texttt{LoanPay}, same loan, same amount, with \texttt{tfLoanLatePayment}.}
\ev{EFAD383F9B622357CE67CBB0518D9F9F5FC27BE611D43F26FD69427713FB608F\\
overpayment 9A7589384C5142703AFFACBB91B12E223EC734F7ACBE51272A87525FE4EE9CA8}
\end{shotblock}

% ───────────── 3 ─────────────
\hh{3 · The first-loss capital covered 0.5\,\% of a defaulted loan}

Both brokers used \tm{CoverRateMinimum} 10\,\% and \tm{CoverRateLiquidation} 5\,\%.

{\setlength{\tabcolsep}{2.8pt}\fontsize{6.8}{8.2}\selectfont
\begin{tabularx}{\linewidth}{@{}r r r r r >{\raggedleft\arraybackslash}X@{}}
\rowcolor{card}
{\bfseries\color{navy}Vault} & {\bfseries\color{navy}Loan} & {\bfseries\color{navy}Cover} & {\bfseries\color{navy}Taken} & {\bfseries\color{navy}Loss} & {\bfseries\color{navy}Share value}\\
\midrule
200 XRP & 50 & 50 & 0.25 & 49.75 & 1.000\,$\rightarrow$\,0.751\\
\rowcolor{card}
5 XRP & 4 & 1 & 0.02 & 3.98 & 1.000\,$\rightarrow$\,0.204\\
\end{tabularx}}

On both vaults the cover taken equals \tm{debt $\times$ 10\% $\times$ 5\%} to the drop. The
xrpl.org Lending Protocol page gives this formula, with an example that covers 1\,\% of the
debt. \tm{CoverRateLiquidation} applies to the minimum cover (\tm{DebtTotal $\times$
CoverRateMinimum}), not to the defaulted amount. After a default, \tm{DebtTotal} is zero and so
is the cover floor.

Second vault, in ledger order:\par\nobreak

{\setlength{\tabcolsep}{2.8pt}\fontsize{6.8}{8.2}\selectfont
\begin{tabularx}{\linewidth}{@{}l >{\raggedright\arraybackslash}X l@{}}
\rowcolor{card}
{\bfseries\color{navy}UTC} & {\bfseries\color{navy}Transaction} & {\bfseries\color{navy}Result}\\
\midrule
14:23:50 & \tm{VaultWithdraw}, all shares & \tm{tecINSUFFICIENT\_FUNDS}\\
\rowcolor{card}
14:27:01 & \tm{LoanManage} \tm{tfLoanDefault} & \tm{tesSUCCESS}\\
14:27:11 & \tm{VaultWithdraw}, all shares, 1.02 XRP & \tm{tesSUCCESS}\\
\rowcolor{card}
14:27:22 & \tm{LoanBrokerCoverWithdraw}, 0.98 XRP & \tm{tesSUCCESS}\\
\end{tabularx}}
\ev{D91BA47E816FDA60F9BBAD7F100B383099D87021B08DE065E841B9E8E34AFE2D\\
3338BA629F8772CDEC07B6F714F2AF3A1ADC8E93642429A54B95CF372FDB2880\\
first vault default 91F458E3E6244E1C3005345C5C83CDFAF0E8A1AF6E5587E100C302CA56743BC7}

\begin{shotblock}
\begin{screen}
\shot{f3a-badge}\shot{f3a-flags}\shot{f3a-meta}
\end{screen}
\capt{Explorer. \texttt{LoanManage} \texttt{tfLoanDefault} moves 0.02 XRP from the broker's pseudo-account (\texttt{rbhF8…}) to the vault's (\texttt{rsR4Q…}).}
\ev{923F7D616B094C82197506CF2867907A93067041E3990D24BD07A19035F73EAC}
\end{shotblock}
\begin{shotblock}
\begin{screen}
\shot{f3b-badge}\shot{f3b-amount}\shot{f3b-date}
\end{screen}
\capt{Explorer. \texttt{LoanBrokerCoverWithdraw} of the remaining cover, 21 s after the default.}
\ev{68DD97D949602FD470E0FA80B2526C3FDCFA3449095F7A7879F4D4430845D45E}
\end{shotblock}

% ───────────── 4 ─────────────
\hh{4 · Share value and \tm{LossUnrealized}}

XLS-65 computes the assets returned on redemption as
\tm{shares $\times$ (AssetsTotal − LossUnrealized) / total shares}, and waives the
\tm{LossUnrealized} deduction when the withdrawer is the sole outstanding shareholder. The
xrpl.org Single Asset Vault page gives the formula and does not mention the exception.

\begin{itemize}
\item The \tm{Vault} node has no \tm{LossUnrealized} field until an impairment writes one. The impairment below moved it from absent to \tm{"3000001"}, with \tm{AssetsTotal} at \tm{"5000002"}.
\item \tm{AssetsAvailable}, \tm{AssetsMaximum}, \tm{PaymentRemaining} and \tm{TotalValueOutstanding} are also absent from their node when zero.
\item The explorer's \emph{Detailed} tab shows ``\emph{It modified a node with type Vault}'' without field values. \tm{LossUnrealized} appears in the \emph{Raw} tab once the node is expanded.
\end{itemize}

\begin{shotblock}
\begin{screen}
\shot{f2-badge}\shot{f2-meta}\vspace{3pt}\shot{f2-raw}
\end{screen}
\capt{Explorer. \texttt{LoanManage} \texttt{tfLoanImpair}: \emph{Detailed} tab, then \emph{Raw} tab with the \texttt{Vault} node expanded.}
\ev{C8678C1C0A10BE889124F40C03017038654B278E2D464FAC54784F55B35318FC}
\end{shotblock}

% ───────────── 5 ─────────────
\hh{5 · Read commands}

\tm{vault\_info} is a \tm{rippled} command. \tm{loan\_info} and \tm{loan\_broker\_info} return
\tm{unknownCmd} and have no documentation page. xrpl.org documents \tm{mpt\_holders} as
Clio-only. The devnet endpoints serve \tm{rippled}, where \tm{mpt\_holders} returns
\tm{unknownCmd}. A depositor's view takes five calls:\par\nobreak

{\setlength{\tabcolsep}{2.8pt}\fontsize{6.8}{8.2}\selectfont
\begin{tabularx}{\linewidth}{@{}l >{\raggedright\arraybackslash}X@{}}
\rowcolor{card}
{\bfseries\color{navy}Call} & {\bfseries\color{navy}Returns}\\
\midrule
\tm{vault\_info} & vault totals, share issuance and share flags\\
\rowcolor{card}
\tm{account\_objects} \tm{mptoken} & the depositor's shares\\
\tm{account\_objects} on the owner & the broker, filtered on \tm{VaultID}\\
\rowcolor{card}
\tm{account\_objects} \tm{loan} & the loans, on the broker's pseudo-account\\
\tm{mpt\_holders} & \tm{unknownCmd}\\
\end{tabularx}}

Six values are computed client-side: gross and net share value, position value, withdrawable
amount, utilisation and broker debt capacity.

% ───────────── 6 ─────────────
\hh{6 · Four result codes cover situations that call for different actions}

{\setlength{\tabcolsep}{2.6pt}\fontsize{6.7}{8}\selectfont
\begin{tabularx}{\linewidth}{@{}l >{\raggedright\arraybackslash}X r@{}}
\rowcolor{card}
{\bfseries\color{navy}Code} & {\bfseries\color{navy}Situation} & {\bfseries\color{navy}Hash}\\
\midrule
\tm{tecINSUFFICIENT\_FUNDS} & \tm{LoanSet} above the cover limit & \hsh{5DC3A6E0}{4147618B}\\
 & \tm{LoanSet} above the vault liquidity & \hsh{B92A90F5}{259F9CAA}\\
 & \tm{LoanBrokerCoverWithdraw} below the floor & \hsh{4FB4A2B0}{262F1370}\\
\rowcolor{card}
\tm{tecLIMIT\_EXCEEDED} & \tm{VaultDeposit} above \tm{AssetsMaximum} & \hsh{132BAF7C}{027B5B4F}\\
\rowcolor{card}
 & \tm{AssetsMaximum} set below \tm{AssetsTotal} & \hsh{37053F15}{A88C3DC6}\\
\rowcolor{card}
 & \tm{DebtMaximum} set below \tm{DebtTotal} & \hsh{82C14ED4}{748E4121}\\
\tm{tecNO\_PERMISSION} & \tm{tfLoanImpair} sent by the borrower & \hsh{F0304194}{5DE2C112}\\
 & \tm{tfLoanUnimpair} on a loan that is not impaired & \hsh{13975FD3}{6AA24537}\\
\rowcolor{card}
\tm{tecNO\_AUTH} & \tm{Payment} of shares to an account without \tm{MPToken} & \hsh{789AAF7E}{5FD512A4}\\
\rowcolor{card}
 & \tm{Payment} of non-transferable shares & \hsh{E9DE504D}{B8BE5FD9}\\
\end{tabularx}}

\tm{tecHAS\_OBLIGATIONS} names a single cause on \tm{VaultDelete}, \tm{LoanBrokerDelete} and
\tm{LoanDelete}.

\begin{shotblock}
\begin{screen}
\shot{f6a-badge}\shot{f6a-status}
\end{screen}
\capt{Explorer. \texttt{LoanSet} above the cover limit.}
\ev{5DC3A6E00F31DA82C69AB77669C64691BD71CAC49C169736F53D7C724147618B}
\end{shotblock}
\begin{shotblock}
\begin{screen}
\shot{f6b-badge}\shot{f6b-status}
\end{screen}
\capt{Explorer. \texttt{LoanSet} above the vault liquidity. Same result code.}
\ev{B92A90F59EBF79108B5414F61A12656F321EBC264E1BB186FCC14CAF259F9CAA}
\end{shotblock}

% ───────────── 7 ─────────────
\hh{7 · The hackathon build reverts the V1.1 \tm{LoanBrokerSet} restriction}

The V1.1 documentation states ``\emph{\texttt{LoanBrokerSet} is restricted on open-ended
vaults}''. The \tm{feature} command reports \tm{LendingProtocolV1\_1} as enabled.
\tm{LoanBrokerSet} returned \tm{tesSUCCESS} on an open-ended vault and on a closed-ended vault.
rippled PR \#8076 added the restriction. The commit \tm{Revert "fix: Reject open-ended vaults at
LoanBrokerSet (\#8076)"} removed it from the lending hackathon branch on 10 September 2026. The other V1.1 rules tested (new fields, date validation, phase locks) are
enforced. The event brief does not mention the revert.
\ev{open-ended 781F54B5E77A768F4C02A54E3C55902AA9F1AC96AA04037AB97CE93FFB5DBDE4\\
closed-ended 1961BE21CC4011F50AA926E1494E9D51F94FC57B0E9C0C1B8F8ADF6FEDD09D40}

% ───────────── 8 ─────────────
\hh{8 · Ports 51233 and 51234}

The documented endpoints use ports 51233 (WebSocket) and 51234 (JSON-RPC). No endpoint is
documented on port 443. On a Wi-Fi network, \tm{nc -z} succeeded on both ports while \tm{curl}
on 51234 returned 0 bytes, and \tm{portquiz.net:51234} timed out. The same signature appeared
five times during the event: three occurrences were traced to the local network, two could not
be attributed. The first cost 45 minutes. The \tm{xrpl.js} \tm{Client} uses WebSocket only. The
faucet ignores the \tm{destination} field and creates a new account.

% ───────────── 9 ─────────────
\hh{9 · Closing a vault to deposits}

The xrpl.org \tm{Vault} entry and XLS-65 define \tm{AssetsMaximum: 0} as no cap. The
\tm{VaultSet} reference does not repeat it. \tm{VaultSet} exposes \tm{Data}, \tm{AssetsMaximum}
and \tm{DomainID}, and no flag closes a vault to deposits. Setting \tm{AssetsMaximum} equal to
\tm{AssetsTotal} blocks deposits: a 1 XRP deposit then returns \tm{tecLIMIT\_EXCEEDED}.
\ev{132BAF7C1B8208E695A453C0EE061BC5E3528FC97571589D8A2EA66F027B5B4F}

\begin{shotblock}
\begin{screen}
\shot{f5a-badge}\shot{f5a-max}\shot{f5a-date}
\end{screen}
\capt{Explorer. \texttt{VaultSet} with \texttt{AssetsMaximum} 0 on a vault capped at 5 XRP and holding 3.}
\ev{A336D71E738AA27C2B449DE05A0FE5C8945733A53420B9B4D8390172404CC1E2}
\end{shotblock}
\begin{shotblock}
\begin{screen}
\shot{f5b-badge}\shot{f5b-amount}\shot{f5b-date}
\end{screen}
\capt{Explorer. Two seconds later, a 4 XRP \texttt{VaultDeposit} brings the same vault to 7 XRP.}
\ev{D27E21CBB4C55F70963BEDB00DFDB28A4C25B8A355619CAE85566CA63E25B956}
\end{shotblock}

% ───────────── 10 ─────────────
\hh{10 · Vault shares carry \tm{lsfMPTCanTrade}}

\tm{vault\_info} returns the share issuance with \tm{Flags: 56}, which is
\tm{lsfMPTCanEscrow} (0x08), \tm{lsfMPTCanTrade} (0x10) and \tm{lsfMPTCanTransfer} (0x20).
\tm{VaultCreate} has no field that sets these flags. The xrpl.org \tm{MPTokenIssuance} reference
describes \tm{lsfMPTCanTrade} as allowing holders to ``\emph{trade their balances using the XRP
Ledger DEX or AMM}''. On the devnet, an \tm{OfferCreate} with a share amount returned
\tm{temDISABLED}. The \tm{feature} list contains \tm{MPTokensV1} and \tm{DynamicMPT}, and no
amendment for MPT trading.

% ───────────── fin ─────────────
\hh{What we ran}

{\setlength{\tabcolsep}{2.6pt}\fontsize{6.7}{8}\selectfont
\begin{tabularx}{\linewidth}{@{}>{\raggedright\arraybackslash}X l r@{}}
\rowcolor{card}
{\bfseries\color{navy}Step} & {\bfseries\color{navy}Transaction} & {\bfseries\color{navy}Hash}\\
\midrule
Open-ended vault & \tm{VaultCreate} & \hsh{A14CB64F}{159502E4}\\
\rowcolor{card}
Deposit, 25 XRP & \tm{VaultDeposit} & \hsh{98A6A315}{FB772A0A}\\
Loan broker & \tm{LoanBrokerSet} & \hsh{A4A619AA}{6480C5A5}\\
\rowcolor{card}
First-loss capital, 5 XRP & \tm{LoanBrokerCoverDeposit} & \hsh{9DEA76E2}{30BF79BB}\\
Origination and drawdown & \tm{LoanSet} & \hsh{7E3006E8}{5F63BAB0}\\
\rowcolor{card}
Repayment & \tm{LoanPay} & \hsh{41B2EAAE}{0E8EDC8C}\\
Withdrawal with yield & \tm{VaultWithdraw} & \hsh{53133F76}{642895BE}\\
\end{tabularx}}

All seven returned \tm{tesSUCCESS}. All 15 XLS-65/66 transaction types were submitted at least
once, with a full \tm{tfLoanImpair} then \tm{tfLoanDefault} cycle.

\hh{Guardrails triggered on purpose}
{\setlength{\tabcolsep}{2.6pt}\fontsize{6.7}{8}\selectfont
\begin{tabularx}{\linewidth}{@{}>{\raggedright\arraybackslash}X l r@{}}
\rowcolor{card}
{\bfseries\color{navy}Transaction} & {\bfseries\color{navy}Result} & {\bfseries\color{navy}Hash}\\
\midrule
\tm{LoanSet} of 5000 XRP, vault holding 25 & \tm{tecINSUFFICIENT\_FUNDS} & \hsh{B92A90F5}{259F9CAA}\\
\rowcolor{card}
\tm{LoanPay} of 1 XRP, instalment 25.5 XRP & \tm{tecINSUFFICIENT\_PAYMENT} & \hsh{77F708D6}{1E0541B4}\\
\tm{LoanBrokerCoverWithdraw} below \tm{CoverRateMinimum} & \tm{tecINSUFFICIENT\_FUNDS} & \hsh{D91DEFCE}{2595A60F}\\
\rowcolor{card}
\tm{LoanBrokerSet} by a non-owner of the vault & \tm{tecNO\_PERMISSION} & \hsh{33742272}{A0FA40A8}\\
\tm{VaultWithdraw} above unlent liquidity & \tm{tecINSUFFICIENT\_FUNDS} & \hsh{FB47313D}{D76DAAC4}\\
\rowcolor{card}
Cover withdrawal above the floor (control) & \tm{tesSUCCESS} & \hsh{90BF019F}{8C4279B2}\\
\end{tabularx}}

\hh{Closed-ended vault, one transaction per phase}
{\setlength{\tabcolsep}{2.6pt}\fontsize{6.7}{8}\selectfont
\begin{tabularx}{\linewidth}{@{}l >{\raggedright\arraybackslash}X l r@{}}
\rowcolor{card}
{\bfseries\color{navy}Phase} & {\bfseries\color{navy}Transaction} & {\bfseries\color{navy}Result} & {\bfseries\color{navy}Hash}\\
\midrule
Subscription & \tm{VaultDeposit} 20 XRP & \tm{tesSUCCESS} & \hsh{2F8B4C3A}{3E25E149}\\
\rowcolor{card}
Investment & \tm{VaultDeposit} 10 XRP & \tm{tecEXPIRED} & \hsh{85D62B25}{B42C92CA}\\
Investment & \tm{VaultWithdraw} 5 XRP & \tm{tecTOO\_SOON} & \hsh{C5449EC0}{60386C5C}\\
\rowcolor{card}
Redemption & \tm{VaultWithdraw} & \tm{tesSUCCESS} & \hsh{F95FA55F}{23F8D8E2}\\
\end{tabularx}}

\hh{Secondary sale of vault shares with TokenEscrow}
{\setlength{\tabcolsep}{2.6pt}\fontsize{6.7}{8}\selectfont
\begin{tabularx}{\linewidth}{@{}>{\raggedright\arraybackslash}X l r@{}}
\rowcolor{card}
{\bfseries\color{navy}Step} & {\bfseries\color{navy}Transaction} & {\bfseries\color{navy}Hash}\\
\midrule
Buyer opts in to the share MPT & \tm{MPTokenAuthorize} & \hsh{79F48365}{B32C5E31}\\
\rowcolor{card}
Buyer locks the payment & \tm{EscrowCreate} & \hsh{1F29F998}{E539FBF8}\\
Seller locks the shares & \tm{EscrowCreate} & \hsh{795D6AFE}{FB204099}\\
\rowcolor{card}
Buyer claims the shares with the preimage & \tm{EscrowFinish} & \hsh{029559F0}{1D26A179}\\
Seller claims the payment with the same preimage & \tm{EscrowFinish} & \hsh{5F359ADB}{50BEF991}\\
\end{tabularx}}

\hh{Lending-related amendments enabled}
\tm{LendingProtocol}, \tm{LendingProtocolV1\_1}, \tm{SingleAssetVault}, \tm{MPTokensV1},
\tm{DynamicMPT}, \tm{TokenEscrow}, \tm{fixTokenEscrowV1}, \tm{fixCleanup3\_2\_0},
\tm{fixCleanup3\_4\_0}, \tm{BatchV1\_1}, \tm{PermissionedDomains}, \tm{Credentials}.

\hh{What worked}
\begin{itemize}
\item \tm{tecHAS\_OBLIGATIONS} names its cause on the three delete transactions.
\item \tm{VaultDelete} removes third-party \tm{MPToken} objects and refunds their reserve.
\item The cover floor is enforced to the drop. \tm{CoverRateMinimum}, \tm{CoverRateLiquidation} and \tm{ManagementFeeRate} return \tm{temINVALID} when changed on an existing broker.
\item The closed-ended vault phases behave as XLS-65 describes: deposits accepted then \tm{tecEXPIRED}, withdrawals \tm{tecTOO\_SOON} then accepted.
\end{itemize}

\hh{Limits}
XRP only, so neither clawback transaction was tested with an issued asset. Payment schedules of
60 to 90 seconds. One build: the hackathon branch of \tm{rippled} 3.4.0-rc1.

\hh{Material in the repo}
\begin{itemize}
\item \tm{scripts/raw-submit.mjs}: a \tm{LoanSet} counterparty signer using \tm{encodeForSigningCounterparty}.
\item \tm{scripts/\_probe-audit-hashes.mjs}: checks every hash cited in a document against the ledger, and separates transaction hashes from ledger object IDs.
\item \tm{deck/capture-explorer.mjs}: the explorer screenshots in this PDF.
\item Observed in passing: the Lending Protocol concept page spells \tm{PrincipleOutstanding} and \tm{depostitor}. \tm{validateVaultCreate} in \tm{xrpl@4.6.0} accepts any number for \tm{WithdrawalPolicy}; the values 0, 2, 3, 99 and 255 reach the network and return \tm{temMALFORMED}.
\end{itemize}

\end{multicols}
\end{document}
